\documentclass[11pt,a4paper]{article}

\usepackage{amsmath,amssymb}
\usepackage{graphicx}
\usepackage{booktabs}
\usepackage{hyperref}
\usepackage[margin=2.5cm]{geometry}
\usepackage{caption}
\usepackage{subcaption}
\usepackage{placeins}
\usepackage{array}
\hypersetup{hidelinks}
\emergencystretch=1em

\title{Mutual Visual HIL Model Comparison for\\Firefly-Inspired Synchronisation}
\author{MSc Student -- BIOE70025}
\date{\today}

\newcommand{\figdir}{figures}
\newcommand{\tabledir}{tables}

\begin{document}
\maketitle

\begin{abstract}
This report analyses a formal mutual hardware-in-the-loop (HIL)
comparison between EAPF Consensus and Kuramoto for firefly-inspired visual
synchronisation. The experiment uses one browser-rendered virtual agent and
one Raspberry Pi visual agent in a bidirectional loop: the Pi camera observes
the virtual flash, the Pi oscillator flashes its GPIO LED, and Pi flash events
are posted back to the frontend. The formal data comprise 18 completed trials:
two models, three Pi initial frequencies, and three repeats with controlled
random initial phases. Under this camera-mediated event-based mutual HIL
condition, EAPF Consensus achieves lower final-window frequency error and much
lower frequency variability than Kuramoto with locked $K=5.0$.
An appendix Kuramoto gain sensitivity sweep over $K \in
\{2.5, 3.0, 3.5, 4.0, 4.5\}$ finds that $K=2.5$ improves behaviour but does
not fully resolve the oscillatory response. EAPF remains the model selected
for the next 2-virtual + 1-Pi mutual HIL stage.
\end{abstract}

\FloatBarrier

% ======================================================================
\section{Introduction}
% ======================================================================

The broader project investigates decentralised visual flashing as a mechanism
for coordinating multi-drone behaviour. Earlier stages established simulation
models, camera flash detection, GPIO LED output, a browser-based visual leader,
a fixed-leader visual HIL comparison, and a simulation-based multi-neighbour
model selection. The fixed-leader setting is an important but simpler one-way
tracking problem: the virtual display acts as a stable external reference and
the Pi follows it without the leader reacting to the follower.

This report extends the evaluation to mutual visual HIL. In this mode, the
virtual agent is no longer a fixed leader. Instead, the Pi observes the virtual
flash through the camera, while its own flashes are sent back to the frontend
over the HTTP API. This creates a closed-loop bidirectional coupling problem.
The purpose of this report is to evaluate which of the two model candidates,
EAPF Consensus and Kuramoto, is more robust in this mutual event-based visual
loop.

This report does not introduce new fixed-leader experiments and does not claim
universal model superiority. It asks a narrower question: based on the formal
mutual visual HIL data currently available, which model is the better
engineering candidate for the next 2-virtual + 1-Pi mixed-reality multi-agent
stage?

% ======================================================================
\section{Experimental Objective}
% ======================================================================

The objective was to compare EAPF Consensus and Kuramoto under matched
1-virtual + 1-Pi mutual HIL conditions. The virtual reference frequency was
2.0 Hz. The Pi oscillator was initialised at 1.2 Hz, 1.5 Hz, or 2.5 Hz to
test convergence across below-reference, moderate-mismatch, and above-reference
conditions. Each model-condition pair was repeated three times with controlled
random initial phases.

The primary outcome is mutual frequency agreement in the final 5 s of a
30 s trial. Secondary outcomes include virtual flash detection reliability,
flash counts, runtime diagnostics, and variability across repeats.

% ======================================================================
\section{Method}
% ======================================================================

\subsection{Mutual HIL setup}

The 1-virtual + 1-Pi mutual HIL loop contains one browser-rendered virtual
agent and one real Pi agent. The frontend displays the virtual flash as a
high-contrast target. The Pi camera detects the virtual flash, updates its
local oscillator, and drives the GPIO LED. When the Pi oscillator fires, the
runner posts a timestamped event to the frontend API. The frontend server then
updates the virtual agent from those Pi flash events.

This bidirectional loop is substantially different from fixed-leader visual
HIL. In fixed-leader mode the visual target is an exogenous, stable reference:
the leader flashes at a fixed frequency and does not react to the follower.
The follower's coupling corrections affect only its own phase and frequency.
In mutual HIL, both agents influence each other, so unstable feedback or
event-timing sensitivity can appear even if a model worked well in the one-way
tracking case.

\subsection{Continuous phase coupling versus binary flash events}

An important architectural distinction affects the model comparison. Kuramoto
is naturally a continuous phase-coupled model: in its standard form the
coupling term $\sin(\theta_j - \theta_i)$ uses continuously available
neighbour phase. In pure simulation, that is a natural assumption.

In the implemented mutual visual HIL system, however, the camera stream is
continuous in time but the binary flash detection signal mainly provides sparse
event timestamps. Continuous neighbour phase is reconstructed from those flash
timestamps---estimated as $\hat\theta = 2\pi (t - t_{\mathrm{last}}) /
T_{\mathrm{est}}$, where $T_{\mathrm{est}}$ is the median of recent inter-flash
intervals. The resulting interaction is therefore an event-derived
approximation to continuous phase coupling.

EAPF Consensus was designed from the start for event-based phase and frequency
tracking. Its corrections are filtered (low-pass with $\alpha=0.2$) and
rate-limited (maximum phase step 0.2 rad, maximum frequency step 0.05 Hz).
Those safeguards are well matched to camera-derived flash timestamps, which
carry inherent jitter from detection latency variability, frame-rate
quantisation, and API communication delay.

This structural mismatch---continuous-phase coupling law operating on
reconstructed phase from sparse events---is central to understanding the
observed behavioural differences between the two models.

\subsection{Formal dataset}

The formal data used here are from the completed mutual visual HIL comparison
run of 2026-06-21. The dataset path is recorded in the reproducibility
appendix (Appendix~\ref{sec:reproducibility}). The dataset includes:

\begin{itemize}
  \item \texttt{aggregate\_metrics.csv}
  \item \texttt{summary\_by\_model\_condition.csv}
  \item \texttt{batch\_metadata.json}
  \item \texttt{run\_metadata.json}
  \item \texttt{chunk\_metadata.json}
  \item per-trial \texttt{metrics\_summary.json}, oscillator logs, detection
        logs, API logs, and flash-event logs
\end{itemize}

\subsection{Models and locked parameters}

The locked model parameters were read from the formal metadata and are shown
in Table~\ref{tab:parameters}. EAPF Consensus uses the event-based
phase/frequency consensus gains and limits selected in the earlier
simulation-based model selection. Kuramoto uses the locked mutual-comparison
gain $K=5.0$.

\begin{table}[htbp]
\centering
\footnotesize
\caption{Locked model parameters used in the formal mutual HIL comparison.}
\label{tab:parameters}
\begin{tabular}{@{}lp{10.5cm}@{}}
\toprule
\textbf{Model} & \textbf{Parameters} \\
\midrule
EAPF Consensus & $g_p=0.02$, $g_f=0.02$, $\alpha_p=\alpha_f=0.2$, $\Delta\theta_{\mathrm{max}}=0.2$ rad, $\Delta f_{\mathrm{max}}=0.05$ Hz \\
Kuramoto & $K=5.0$ \\
\bottomrule
\end{tabular}
\end{table}


\subsection{Trial schedule and seed control}

The formal run was executed as three chunks, one per Pi initial frequency. The
chunk metadata records all three completed chunks (Table~\ref{tab:chunks}).
Each chunk completed six trials: two models by three repeats. Across all
chunks this gives 18 completed trials and zero failed trials.

Random initial phases were enabled. The base seed was 20260620. Trial seeds
were assigned using a canonical schedule over Pi frequencies 1.2, 1.5, and
2.5 Hz, so chunking did not change the phase schedule.

\begin{table}[htbp]
\centering
\footnotesize
\caption{Formal chunk validation records from chunk metadata.}
\label{tab:chunks}
\begin{tabular}{@{}llrrl@{}}
\toprule
\textbf{Chunk} & \textbf{Freq pairs} & \textbf{Completed} & \textbf{Failed} & \textbf{Started} \\
\midrule
pi\_1p2 & 2:1.2 & 6 & 0 & 2026-06-21T19:41:38.144053 \\
pi\_1p5 & 2:1.5 & 6 & 0 & 2026-06-21T19:48:33.342194 \\
pi\_2p5 & 2:2.5 & 6 & 0 & 2026-06-21T19:54:48.310484 \\
\bottomrule
\end{tabular}
\end{table}


\subsection{Metrics}

The core performance metrics are:

\begin{itemize}
  \item actual virtual flash count ratio (FCR), using actual virtual flashes
        as the denominator;
  \item final 5 s mean absolute frequency error between Pi and virtual agents;
  \item absolute difference between the final 5 s mean Pi and virtual
        frequencies;
  \item final 5 s mean virtual and Pi frequencies;
  \item final-window frequency variability across each trial;
  \item Pi flash count, detected virtual flash count, and actual virtual flash
        count.
\end{itemize}

Runtime diagnostics were derived from per-trial logs. Capture FPS is estimated
from detection-log sample count divided by trial duration. Processing FPS is
the reciprocal of the mean camera processing time. API POST latency is measured
from \texttt{api\_events.csv}.

% ======================================================================
\section{Results}
% ======================================================================

\subsection{Dataset validation}

All expected formal trials were present: two models, three Pi initial
frequencies, and three repeats. Each trial had a per-trial
\texttt{metrics\_summary.json}; no formal trial was excluded for corruption or
incompleteness. The full validation table is included in
Appendix~\ref{sec:appendix_validation}.

\begin{table}[htbp]
\centering
\footnotesize
\caption{Mutual HIL overall results across all formal trials (1 virtual + 1 Pi, 30 s, 3 repeats).}
\label{tab:overall}
\begin{tabular}{@{}lrrrrrr@{}}
\toprule
\textbf{Model} & \textbf{Trials} & \textbf{Actual FCR} & \textbf{5 s MAE} & \textbf{Mean diff.} & \textbf{Virt. std} & \textbf{Pi std} \\
\midrule
EAPF Consensus & 9 & 0.946 $\pm$ 0.028 & 0.046 $\pm$ 0.038 & 0.043 & 0.001 & 0.011 \\
Kuramoto & 9 & 0.903 $\pm$ 0.071 & 0.413 $\pm$ 0.095 & 0.199 & 0.247 & 0.393 \\
\bottomrule
\end{tabular}
\end{table}


Table~\ref{tab:overall} shows the main cross-condition result. EAPF Consensus
has a mean final 5 s frequency error of 0.046 Hz, compared with 0.413 Hz for
Kuramoto (locked $K=5.0$). The final-window frequency variability is also much
lower for EAPF: mean virtual frequency standard deviation is 0.001 Hz for EAPF
and 0.247 Hz for Kuramoto. The Pi-side final-window standard deviation shows
the same pattern (0.011 Hz vs.\ 0.393 Hz).

\begin{table}[htbp]
\centering
\footnotesize
\caption{Mean performance by model and Pi initial frequency condition. Frequency metrics are in Hz.}
\label{tab:condition_summary}
\begin{tabular}{@{}llrrrrrr@{}}
\toprule
\textbf{Pi} & \textbf{Model} & \textbf{Trials} & \textbf{FCR} & \textbf{5 s MAE} & \textbf{Mean diff.} & \textbf{Virt. final} & \textbf{Pi final} \\
\midrule
1.2 Hz & EAPF & 3 & 0.960 $\pm$ 0.025 & 0.029 $\pm$ 0.019 & 0.029 & 1.908 & 1.879 \\
1.2 Hz & Kuramoto & 3 & 0.938 $\pm$ 0.057 & 0.440 $\pm$ 0.084 & 0.371 & 1.930 & 1.560 \\
1.5 Hz & EAPF & 3 & 0.937 $\pm$ 0.043 & 0.084 $\pm$ 0.045 & 0.084 & 1.949 & 1.866 \\
1.5 Hz & Kuramoto & 3 & 0.946 $\pm$ 0.037 & 0.348 $\pm$ 0.133 & 0.119 & 1.848 & 1.745 \\
2.5 Hz & EAPF & 3 & 0.941 $\pm$ 0.010 & 0.026 $\pm$ 0.015 & 0.015 & 2.126 & 2.139 \\
2.5 Hz & Kuramoto & 3 & 0.824 $\pm$ 0.043 & 0.451 $\pm$ 0.040 & 0.107 & 2.193 & 2.086 \\
\bottomrule
\end{tabular}
\end{table}


\subsection{Robustness and detection}

Figure~\ref{fig:overview} summarises two key axes: virtual flash detection
reliability and final 5 s frequency agreement. Bars show the mean across three
repeats per model-condition pair; black points show individual trial values.
Both models generally retain usable virtual flash detection (most trials above
the indicative FCR threshold of 0.85), but EAPF is more consistent. The
difference is larger for frequency agreement, where Kuramoto has much higher
error across all Pi initial frequency conditions.

\begin{figure}[htbp]
\centering
\includegraphics[width=0.95\textwidth]{\figdir/fig1_robustness_overview.pdf}
\caption{Mutual HIL robustness overview (1 virtual + 1 Pi, 30 s trials).
Bars show mean across three repeats per model-condition pair; black points
show individual trial values. The dashed line in the left panel marks an
indicative actual-FCR threshold of 0.85.}
\label{fig:overview}
\end{figure}

\subsection{Frequency agreement}

Figure~\ref{fig:final_freqs} shows the final 5 s mean frequencies by model and
condition. EAPF keeps both virtual and Pi final frequencies close to the
2.0 Hz reference. Kuramoto final frequencies are more dispersed and
condition-dependent. Figure~\ref{fig:agreement} confirms this in scatter form:
EAPF points cluster closer to the $y=x$ agreement line.

\begin{figure}[htbp]
\centering
\begin{subfigure}{0.48\textwidth}
  \centering
  \includegraphics[width=\textwidth]{\figdir/fig2_virtual_final_frequency.pdf}
  \caption{Virtual final frequency.}
  \label{fig:virtual_final}
\end{subfigure}
\hfill
\begin{subfigure}{0.48\textwidth}
  \centering
  \includegraphics[width=\textwidth]{\figdir/fig3_pi_final_frequency.pdf}
  \caption{Pi final frequency.}
  \label{fig:pi_final}
\end{subfigure}
\caption{Final 5 s mean frequencies by model and condition. Bars show
model-condition means; black points are individual trials.}
\label{fig:final_freqs}
\end{figure}

\begin{figure}[htbp]
\centering
\includegraphics[width=0.62\textwidth]{\figdir/fig4_frequency_agreement_scatter.pdf}
\caption{Final frequency agreement. Each point is one trial. The dashed line
is the $y=x$ reference. EAPF points cluster closer to the agreement line.}
\label{fig:agreement}
\end{figure}

\subsection{Representative trajectories}

Figure~\ref{fig:timeseries} compares representative high-mismatch trials
under the Pi 2.5 Hz initial condition (repeat 3). EAPF shows comparatively
smooth virtual frequency adaptation. Kuramoto exhibits visibly larger
oscillation in both the virtual and Pi frequency traces. This matches the
observed frontend behaviour: the frontend in feedback-ON mode displays
server \texttt{flash\_on}, which is generated from server-side oscillator
phase crossings.

\begin{figure}[htbp]
\centering
\includegraphics[width=0.95\textwidth]{\figdir/fig6_representative_frequency_timeseries.pdf}
\caption{Representative per-trial frequency trajectories for the Pi 2.5 Hz
initial condition, repeat 3 (30 s trials). The virtual frequency trace comes
from \texttt{/api/agents}; the Pi trace comes from the local oscillator log.}
\label{fig:timeseries}
\end{figure}

\subsection{Flash counts and runtime diagnostics}

Figure~\ref{fig:counts} shows the flash-count diagnostics. Counts are broadly
consistent with the 30 s trial duration and observed final frequencies. EAPF
has slightly higher detected virtual flash counts overall, while Kuramoto
shows lower actual FCR in the Pi 2.5 Hz condition.

\begin{figure}[htbp]
\centering
\includegraphics[width=0.95\textwidth]{\figdir/fig8_flash_counts.pdf}
\caption{Pi flash count, detected virtual flash count, and actual virtual
flash count by model and condition. Bars show means; points show individual
trials.}
\label{fig:counts}
\end{figure}

\begin{table}[htbp]
\centering
\footnotesize
\caption{Runtime and detection diagnostics derived from per-trial logs.}
\label{tab:runtime}
\begin{tabular}{@{}lrrrrrr@{}}
\toprule
\textbf{Model} & \textbf{Cap. FPS} & \textbf{Proc. FPS} & \textbf{POST ms} & \textbf{API OK} & \textbf{Loop Hz} & \textbf{Events} \\
\midrule
EAPF Consensus & 21.8 $\pm$ 1.0 & 26.3 $\pm$ 0.8 & 18.1 $\pm$ 7.7 & 1.000 & 29.99 & 56.2 \\
Kuramoto & 22.1 $\pm$ 0.8 & 26.6 $\pm$ 0.9 & 21.9 $\pm$ 10.5 & 1.000 & 29.97 & 52.1 \\
\bottomrule
\end{tabular}
\end{table}


Runtime diagnostics in Table~\ref{tab:runtime} do not show a model-specific
runtime failure. API success ratio was 1.000 for both models. Capture and
processing rates were similar. This supports the interpretation that the main
model difference is not an API or camera throughput artifact.

\begin{figure}[htbp]
\centering
\includegraphics[width=0.95\textwidth]{\figdir/fig7_runtime_diagnostics.pdf}
\caption{Runtime diagnostics derived from detection and API logs. Boxes show
per-trial distributions across the 9 trials per model.}
\label{fig:runtime}
\end{figure}

\subsection{Initial phase sensitivity}

The formal batch used random initial phases with controlled seeds. Figure
\ref{fig:phase} plots initial phase difference against final 5 s frequency
error. With only three repeats per condition, this is an exploratory
diagnostic rather than a formal phase-sensitivity study. It does show that
Kuramoto's larger errors are not confined to a single initial phase offset.

\begin{figure}[htbp]
\centering
\includegraphics[width=0.78\textwidth]{\figdir/fig5_initial_phase_sensitivity.pdf}
\caption{Initial phase difference versus final 5 s mean absolute frequency
error. Each point is one trial.}
\label{fig:phase}
\end{figure}

\FloatBarrier

% ======================================================================
\section{Discussion}
% ======================================================================

\subsection{EAPF versus Kuramoto in mutual HIL}

The formal mutual visual HIL data indicate that EAPF Consensus is more stable
in the current setup. With locked $K=5.0$, its final-window frequency errors
are lower, its final-window frequency variability is much lower, and its
behaviour is consistent across the three Pi initial frequencies. Kuramoto
remains a useful comparison baseline, but in this mutual loop it shows larger
fluctuations and poorer final-window agreement.

This conclusion is specific to the implemented event-based visual HIL system.
It should not be read as a universal statement that EAPF is always superior to
Kuramoto. Rather, EAPF appears better matched to this project's sensing and
feedback pathway: timestamped flash events, camera-mediated detections, and
explicit event-based phase/frequency correction with rate-limiting and
filtering.

\subsection{Why Kuramoto behaviour differs from fixed-leader results}

The mutual HIL result does not contradict earlier fixed-leader visual results.
Fixed-leader HIL is a one-way tracking task: the virtual leader flashes at a
stable, externally controlled frequency and does not react to the Pi. Feedback
perturbations affect only the follower, and cannot be amplified through the
visual loop. In mutual HIL, both agents influence the timing that the other
observes, forming a bidirectional closed loop.

Kuramoto can track a stable fixed visual reference well because the reference
phase is not perturbed by the follower's own output. In mutual HIL, however,
the virtual agent's phase is updated from Pi flash events, so the virtual
visual stimulus itself carries the oscillations introduced by the Kuramoto
phase-velocity correction. This feedback path does not exist in fixed-leader
mode, which explains why a model can perform well in one-way tracking yet
become less stable in mutual HIL.

The frontend instability observed during Kuramoto trials is consistent with
this interpretation. The virtual Kuramoto update uses a direct instantaneous
phase-velocity correction $\dot\theta = \omega + K \sin(\hat\theta_j -
\theta)$. At the locked value $K=5.0$, that correction can produce substantial
short-timescale frequency swings. The frontend then displays the server's
\texttt{flash\_on} state, so visible flashing inherits those phase-crossing
fluctuations.

EAPF avoids much of this behaviour because its corrections are filtered
(low-pass $\alpha = 0.2$) and rate-limited ($\Delta\theta_{\max} = 0.2$ rad,
$\Delta f_{\max} = 0.05$ Hz). This design is naturally better suited to the
event-derived, latency-affected phase information available from camera flash
detection.

\subsection{Limitations}

The formal sample size is modest: three repeats per model-condition pair
(18 trials total). The interaction is camera-mediated and uses a
screen-rendered target, which is not identical to drone-mounted
LED-to-camera interaction. The Pi and frontend communicate through an HTTP
API, and the browser display uses a fixed flash hold after server-side phase
crossings. Kuramoto's gain was inherited from the earlier simulation
comparison setting ($K=5.0$) and may not be optimal for mutual visual HIL.
A gain sensitivity analysis is presented in
Appendix~\ref{sec:k_sensitivity}.

% ======================================================================
\section{Conclusion}
% ======================================================================

The formal mutual visual HIL comparison (1 virtual + 1 Pi, 30 s trials,
locked $K=5.0$ for Kuramoto) supports carrying EAPF Consensus forward to the
next 2-virtual + 1-Pi mutual HIL stage. EAPF achieved lower final-window
frequency error and substantially lower final-window frequency variability
than Kuramoto across all tested Pi initial frequencies.

Kuramoto is retained as a meaningful baseline. An appendix Kuramoto gain
sensitivity sweep (Appendix~\ref{sec:k_sensitivity}) tested $K \in \{2.5,
3.0, 3.5, 4.0, 4.5\}$ and found that $K=2.5$ yields the best behaviour
according to the appendix ranking score. However, even at $K=2.5$, Kuramoto
does not achieve the final-window stability of EAPF (see
Appendix~\ref{sec:best_k_vs_eapf}). This supports the interpretation that the
limitation is not merely an overly large coupling gain, but also the
adaptation of a continuous-phase coupling law to sparse, event-derived phase
estimates from camera flash detection.

The next experimental direction is therefore an EAPF-based 2-virtual + 1-Pi
mutual HIL test, which will extend the mutual loop to three agents in a
mixed-reality configuration.

% ======================================================================
\appendix
\section{Trial Validation}
\label{sec:appendix_validation}

\begingroup
\scriptsize
\begin{table}[htbp]
\centering
\footnotesize
\caption{Per-trial validation for the formal mutual HIL dataset.}
\label{tab:trial_validation}
\begin{tabular}{@{}p{4.55cm}p{2.1cm}rrlll@{}}
\toprule
\textbf{Trial} & \textbf{Model} & \textbf{Pi Hz} & \textbf{Rep.} & \textbf{Aggregate} & \textbf{Metrics} & \textbf{Status} \\
\midrule
eapf\_\allowbreak{}consensus\_\allowbreak{}V2\_\allowbreak{}P1.2\_\allowbreak{}r01 & EAPF & 1.2 & 1 & yes & yes & included \\
kuramoto\_\allowbreak{}V2\_\allowbreak{}P1.2\_\allowbreak{}r01 & Kuramoto & 1.2 & 1 & yes & yes & included \\
eapf\_\allowbreak{}consensus\_\allowbreak{}V2\_\allowbreak{}P1.5\_\allowbreak{}r01 & EAPF & 1.5 & 1 & yes & yes & included \\
kuramoto\_\allowbreak{}V2\_\allowbreak{}P1.5\_\allowbreak{}r01 & Kuramoto & 1.5 & 1 & yes & yes & included \\
eapf\_\allowbreak{}consensus\_\allowbreak{}V2\_\allowbreak{}P2.5\_\allowbreak{}r01 & EAPF & 2.5 & 1 & yes & yes & included \\
kuramoto\_\allowbreak{}V2\_\allowbreak{}P2.5\_\allowbreak{}r01 & Kuramoto & 2.5 & 1 & yes & yes & included \\
eapf\_\allowbreak{}consensus\_\allowbreak{}V2\_\allowbreak{}P1.2\_\allowbreak{}r02 & EAPF & 1.2 & 2 & yes & yes & included \\
kuramoto\_\allowbreak{}V2\_\allowbreak{}P1.2\_\allowbreak{}r02 & Kuramoto & 1.2 & 2 & yes & yes & included \\
eapf\_\allowbreak{}consensus\_\allowbreak{}V2\_\allowbreak{}P1.5\_\allowbreak{}r02 & EAPF & 1.5 & 2 & yes & yes & included \\
kuramoto\_\allowbreak{}V2\_\allowbreak{}P1.5\_\allowbreak{}r02 & Kuramoto & 1.5 & 2 & yes & yes & included \\
eapf\_\allowbreak{}consensus\_\allowbreak{}V2\_\allowbreak{}P2.5\_\allowbreak{}r02 & EAPF & 2.5 & 2 & yes & yes & included \\
kuramoto\_\allowbreak{}V2\_\allowbreak{}P2.5\_\allowbreak{}r02 & Kuramoto & 2.5 & 2 & yes & yes & included \\
eapf\_\allowbreak{}consensus\_\allowbreak{}V2\_\allowbreak{}P1.2\_\allowbreak{}r03 & EAPF & 1.2 & 3 & yes & yes & included \\
kuramoto\_\allowbreak{}V2\_\allowbreak{}P1.2\_\allowbreak{}r03 & Kuramoto & 1.2 & 3 & yes & yes & included \\
eapf\_\allowbreak{}consensus\_\allowbreak{}V2\_\allowbreak{}P1.5\_\allowbreak{}r03 & EAPF & 1.5 & 3 & yes & yes & included \\
kuramoto\_\allowbreak{}V2\_\allowbreak{}P1.5\_\allowbreak{}r03 & Kuramoto & 1.5 & 3 & yes & yes & included \\
eapf\_\allowbreak{}consensus\_\allowbreak{}V2\_\allowbreak{}P2.5\_\allowbreak{}r03 & EAPF & 2.5 & 3 & yes & yes & included \\
kuramoto\_\allowbreak{}V2\_\allowbreak{}P2.5\_\allowbreak{}r03 & Kuramoto & 2.5 & 3 & yes & yes & included \\
\bottomrule
\end{tabular}
\end{table}

\endgroup

% ======================================================================
\section{Kuramoto Gain Sensitivity}
\label{sec:k_sensitivity}

\subsection{Motivation}

In the main mutual visual HIL comparison (Section~4), Kuramoto with locked
$K=5.0$ showed visible frequency fluctuation and higher final-window
variability than EAPF. This appendix tests whether lower Kuramoto coupling
gains reduce the oscillation.

\subsection{Theory framing}

For a frequency mismatch $\Delta f$, the angular mismatch is $\Delta\omega =
2\pi\,\Delta f$. A one-way driven intuition gives a critical coupling
$K_c \approx \Delta\omega$ for frequency locking. A symmetric two-oscillator
mutual Kuramoto intuition gives $K_c \approx \Delta\omega / 2$, because each
oscillator carries half the coupling load.

This HIL system is not an ideal continuous Kuramoto system, because neighbour
phase is reconstructed from binary flash events rather than observed
continuously. The tested $K$ values are therefore theory-informed diagnostics,
not exact theoretical thresholds.

\subsection{Experimental design}

\begin{itemize}
  \item virtual initial frequency = 2.0 Hz
  \item Pi initial frequencies = 1.2, 1.5, 2.5 Hz (matched to main comparison)
  \item $K$ values = 2.5, 3.0, 3.5, 4.0, 4.5
  \item repeats = 2 per condition
  \item duration = 30 s
  \item random phases enabled; phase schedule paired across $K$ values
  \item chunked execution; three chunks of 10 trials each; zero failed trials
\end{itemize}

The main locked-parameter Kuramoto comparison uses $K=5.0$ and is not replaced
by this appendix. This sweep is a secondary sensitivity analysis.

\subsection{Results}

\begin{table}[htbp]
\centering
\footnotesize
\caption{Kuramoto K-sweep: overall results across all Pi initial frequencies (2 repeats per condition).}
\label{tab:k_sweep_overall}
\begin{tabular}{@{}rrrrrrrr@{}}
\toprule
\textbf{K} & \textbf{Trials} & \textbf{FCR} & \textbf{5 s MAE} & \textbf{Mean diff.} & \textbf{Virt. std} & \textbf{Pi std} & \textbf{Score} \\
\midrule
2.5 & 6 & 0.844 & 0.432 & 0.392 & 0.213 & 0.257 & 0.847 \\
3.0 & 6 & 0.864 & 0.320 & 0.261 & 0.198 & 0.277 & 1.730 \\
3.5 & 6 & 0.873 & 0.362 & 0.249 & 0.262 & 0.286 & 2.498 \\
4.0 & 6 & 0.845 & 0.399 & 0.213 & 0.273 & 0.368 & 2.703 \\
4.5 & 6 & 0.868 & 0.332 & 0.123 & 0.179 & 0.344 & 2.461 \\
\bottomrule
\end{tabular}
\end{table}


Table~\ref{tab:k_sweep_overall} summarises results across all Pi initial
frequencies. $K=2.5$ achieves the best (lowest) appendix ranking score of
0.847. The ranking score penalises final-window mean absolute error, frequency
variability, detection-FCR shortfalls, and diagnostic warning counts. Lower
$K$ values generally produce fewer Kuramoto debug warnings (mean 1.8 at
$K=2.5$ vs.\ 87.8 at $K=4.5$), consistent with reduced aggressive
phase-velocity correction.

\begin{table}[htbp]
\centering
\footnotesize
\caption{Kuramoto K-sweep results by Pi initial frequency condition.}
\label{tab:k_sweep_condition}
\begin{tabular}{@{}lrrrrrrr@{}}
\toprule
\textbf{Pi} & \textbf{K} & \textbf{Trials} & \textbf{FCR} & \textbf{5 s MAE} & \textbf{Virt. std} & \textbf{Pi std} & \textbf{Score} \\
\midrule
1.2 Hz & 2.5 & 2 & 0.877 & 0.687 & 0.263 & 0.251 & 1.029 \\
1.5 Hz & 2.5 & 2 & 0.935 & 0.315 & 0.196 & 0.258 & 0.632 \\
2.5 Hz & 2.5 & 2 & 0.721 & 0.293 & 0.180 & 0.261 & 0.881 \\
1.2 Hz & 3.0 & 2 & 0.952 & 0.560 & 0.355 & 0.324 & 2.140 \\
1.5 Hz & 3.0 & 2 & 0.900 & 0.213 & 0.155 & 0.272 & 1.103 \\
2.5 Hz & 3.0 & 2 & 0.740 & 0.186 & 0.086 & 0.234 & 1.946 \\
1.2 Hz & 3.5 & 2 & 0.932 & 0.618 & 0.464 & 0.360 & 2.819 \\
1.5 Hz & 3.5 & 2 & 0.918 & 0.151 & 0.093 & 0.158 & 0.977 \\
2.5 Hz & 3.5 & 2 & 0.767 & 0.318 & 0.229 & 0.340 & 3.698 \\
1.2 Hz & 4.0 & 2 & 0.934 & 0.476 & 0.350 & 0.362 & 2.433 \\
1.5 Hz & 4.0 & 2 & 0.844 & 0.370 & 0.291 & 0.371 & 2.104 \\
2.5 Hz & 4.0 & 2 & 0.757 & 0.352 & 0.178 & 0.370 & 3.572 \\
1.2 Hz & 4.5 & 2 & 0.931 & 0.315 & 0.063 & 0.389 & 1.911 \\
1.5 Hz & 4.5 & 2 & 0.938 & 0.376 & 0.326 & 0.260 & 2.159 \\
2.5 Hz & 4.5 & 2 & 0.736 & 0.306 & 0.150 & 0.383 & 3.311 \\
\bottomrule
\end{tabular}
\end{table}


\begin{table}[htbp]
\centering
\footnotesize
\caption{Kuramoto K ranking from the appendix sensitivity sweep. Lower score is better.}
\label{tab:k_ranking}
\begin{tabular}{@{}rrrrr@{}}
\toprule
\textbf{Rank} & \textbf{K} & \textbf{Score} & \textbf{5 s MAE} & \textbf{Fan-out warnings} \\
\midrule
1 & 2.5 & 0.847 & 0.432 & 1.8 \\
2 & 3.0 & 1.730 & 0.320 & 52.7 \\
3 & 4.5 & 2.461 & 0.332 & 87.8 \\
4 & 3.5 & 2.498 & 0.362 & 88.5 \\
5 & 4.0 & 2.703 & 0.399 & 92.5 \\
\bottomrule
\end{tabular}
\end{table}


Table~\ref{tab:k_sweep_condition} breaks results down by Pi initial frequency.
$K=2.5$ is the best-performing gain in every condition. Table~\ref{tab:k_ranking}
shows the overall ranking.

\begin{figure}[htbp]
\centering
\includegraphics[width=0.95\textwidth]{\figdir/figA1_k_sweep_final5_mae.pdf}
\caption{Final 5 s mean absolute frequency error vs.\ Kuramoto $K$, by Pi
initial frequency. Filled diamonds show the mean across 2 repeats; open
circles show individual trial values.}
\label{fig:k_sweep_mae}
\end{figure}

\begin{figure}[htbp]
\centering
\includegraphics[width=0.95\textwidth]{\figdir/figA2_k_sweep_frequency_std.pdf}
\caption{Final 5 s frequency standard deviation vs.\ $K$. Top row: virtual
frequency std; bottom row: Pi frequency std.}
\label{fig:k_sweep_std}
\end{figure}

\begin{figure}[htbp]
\centering
\includegraphics[width=0.95\textwidth]{\figdir/figA3_k_sweep_actual_fcr.pdf}
\caption{Actual virtual flash count ratio vs.\ $K$. The dashed line marks
FCR = 0.85.}
\label{fig:k_sweep_fcr}
\end{figure}

\begin{figure}[htbp]
\centering
\includegraphics[width=0.55\textwidth]{\figdir/figA4_k_sweep_ranking_score.pdf}
\caption{Appendix ranking score vs.\ $K$ (lower is better). $K=2.5$ is best.
Score formula: $\mathrm{score} = \mathrm{MAE} + 0.5(\sigma_{\mathrm{virt}} +
\sigma_{\mathrm{Pi}}) + 2\cdot\max(0,0.90-\mathrm{FCR}) +
5\cdot\max(0,\mathrm{FCR}-1.02) + 0.02\cdot(\mathrm{warnings})$,
with zero failed trials.}
\label{fig:k_sweep_ranking}
\end{figure}

\subsection{Interpretation}

Lower $K$ reduces the aggressive instantaneous phase-velocity correction
compared with $K=5.0$. $K=2.5$ is the best tested value according to the
appendix ranking. However, the sweep shows that Kuramoto does not fully
converge to smooth, low-variability behaviour in the mutual visual HIL loop
even at this best gain. The remaining oscillation suggests that the issue is
not only gain magnitude, but also the mismatch between continuous-phase
Kuramoto coupling and event-based visual flash timing.

Reducing the Kuramoto gain improves the behaviour relative to the locked
$K=5.0$ setting, but does not fully remove the oscillatory response. This
supports the interpretation that the limitation is not merely an overly large
coupling gain, but also the adaptation of a continuous-phase coupling law to
sparse camera-derived flash events.

% ======================================================================
\section{Best tested Kuramoto gain versus EAPF}
\label{sec:best_k_vs_eapf}

This section compares the existing main EAPF data with the best tested
Kuramoto gain $K=2.5$ from the appendix sweep. No new trials were run; all
data are from the completed formal datasets.

The comparison should be interpreted with care:
\begin{itemize}
  \item EAPF has 3 repeats per condition in the main dataset.
  \item Kuramoto $K=2.5$ has 2 repeats per condition in the appendix sweep.
  \item The phase schedules are controlled and paired, but sample sizes differ.
  \item This is a secondary appendix comparison, not the main
        locked-parameter result.
\end{itemize}

\begin{table}[htbp]
\centering
\footnotesize
\caption{EAPF Consensus (main) versus best-appendix Kuramoto K=2.5. Values shown as ``EAPF / K=2.5''. Sample sizes differ: EAPF 3 repeats, K=2.5 has 2 repeats.}
\label{tab:best_k_vs_eapf}
\begin{tabular}{@{}lrrrrr@{}}
\toprule
\textbf{Pi} & \textbf{N (E/K)} & \textbf{5 s MAE (Hz)} & \textbf{Virt. std (Hz)} & \textbf{Pi std (Hz)} & \textbf{Actual FCR} \\
\midrule
1.2 Hz & 3 & 2 & 0.029  /  0.687 & 0.000  /  0.263 & 0.008  /  0.251 & 0.960  /  0.877 \\
1.5 Hz & 3 & 2 & 0.084  /  0.315 & 0.001  /  0.196 & 0.003  /  0.258 & 0.937  /  0.935 \\
2.5 Hz & 3 & 2 & 0.026  /  0.293 & 0.001  /  0.180 & 0.022  /  0.261 & 0.941  /  0.721 \\
\bottomrule
\end{tabular}
\end{table}


Table~\ref{tab:best_k_vs_eapf} compares the two models across conditions.
EAPF maintains lower final-window frequency variability (virtual std and Pi
std) in all conditions. The 5 s MAE values are more comparable at the
moderate-mismatch Pi 1.5 Hz condition, but EAPF is consistently lower at the
higher-mismatch Pi 1.2 Hz and Pi 2.5 Hz conditions.

\begin{figure}[htbp]
\centering
\includegraphics[width=0.95\textwidth]{\figdir/figB1_eapf_vs_k25_overview.pdf}
\caption{EAPF Consensus (main, 3 repeats) vs.\ Kuramoto $K=2.5$ (appendix,
2 repeats). Left: actual virtual FCR; right: final 5 s MAE. Bars show means;
points show individual trials.}
\label{fig:best_k_overview}
\end{figure}

\begin{figure}[htbp]
\centering
\includegraphics[width=0.85\textwidth]{\figdir/figB2_eapf_vs_k25_frequency_std.pdf}
\caption{Final-window frequency standard deviation: EAPF vs.\ Kuramoto
$K=2.5$. Left: virtual frequency std; right: Pi frequency std.}
\label{fig:best_k_std}
\end{figure}

\subsection{Interpretation}

$K=2.5$ improves Kuramoto relative to $K=5.0$ on several metrics, especially
in reducing the excessive warning/oscillation score. However, EAPF still has
much lower final-window frequency variability and better stability across
conditions.

Therefore the model selection for the next 2-virtual + 1-Pi mutual HIL test
remains EAPF Consensus. This conclusion does not claim that EAPF universally
outperforms Kuramoto in all possible implementations. It states that EAPF is
better suited to this project's current binary-flash, camera-mediated,
event-based mutual HIL setup.

\section{Reproducibility}
\label{sec:reproducibility}

The main formal dataset (1 virtual + 1 Pi mutual HIL comparison) is stored at:

\begin{quote}
\path{experiments/logs/step5b_formal_chunked/formal_step5b_chunked_20260621}
\end{quote}

The Kuramoto K-sensitivity appendix dataset is stored at:

\begin{quote}
\path{experiments/logs/step5b_kuramoto_k_sensitivity/k_sweep_20260621_v4}
\end{quote}

The figures and tables in this report were generated by:

\begin{quote}
\path{docs/mutual_visual_hil_report/generate_mutual_hil_analysis.py}
\end{quote}

The directory names \texttt{step5b\_formal\_chunked} and
\texttt{step5b\_kuramoto\_k\_sensitivity}
in the experiment log paths are internal folder labels only. The conceptual
naming used throughout this report is ``mutual visual HIL'' and
``Kuramoto gain sensitivity appendix.''

Derived copies of the core CSV/JSON inputs and intermediate summary tables
are stored under:

\begin{quote}
\path{docs/mutual_visual_hil_report/derived/}
\end{quote}

\end{document}
